\section{Q2: Fixed-Anion QUBO}
\label{sec:q2-qubo}

\subsection{Fixed-anion objective and scoring input}

Question~2 fixes the anion group at $X=\mathrm{NO}_2$ and seeks the
ordered cation pair with the lowest predicted violation.  The search domain
is file~07, which contains all ordered pairs from the 52-element cation set
($52^2=2,704$ candidates).  We optimize the frozen candidate-scoring model
defined in Section~\ref{sec:problem-formulation}:
\begin{equation}
  (A_2^*,B_2^*)=
  \arg\min_{A,B\in\mathcal E}
  \widehat V_{23}(A,B,\mathrm{NO}_2).
  \label{eq:q2-fixed-objective}
\end{equation}
Here, $\widehat V_{23}$ denotes the selected \texttt{hgb\_gap\_mix}
specification refitted on files~02 and~03.  The model uses the ordered
composition fields $(A,B,X)$ and the deterministic composition features
defined by the frozen specification.  File~07 scores serve as optimization
inputs, not as additional validation observations; no physical labels from
files~07, 11, or the hidden pool are used in fitting.  We first enumerate
the frozen score table to identify its proxy argmin and runner-up.  The
integer QUBO is then used to test whether this already known argmin can be
encoded and sampled within the recorded coefficient limit.  Because the
group minimum is read from the enumerated score table and inserted into the
cost transform, this full-pool construction is oracle-assisted by design;
its result is an encoding and sampling test, not an independent discovery
of the argmin.

\subsection{One-hot encoding and accepted integer cost}

Let $a_i,b_j\in\{0,1\}$ select the $i$th A-site element and the $j$th
B-site element.  There are 52 variables in each block, hence 104 logical
variables in total.  Legal assignments satisfy the two one-hot constraints
in Equation~\eqref{eq:one-hot-constraints}.  The accepted implementation does not insert the raw floating-point
scores directly into the QUBO.  For a fixed group $X$, it first subtracts
the group minimum
$m_X=\min_{i,j}\widehat V_{23}(A_i,B_j,X)$ and forms
\begin{equation}
 c_{ij}^{(X)} = 4\,\operatorname{nint}\!\left(
   \frac{\min\{\widehat V_{23}(A_i,B_j,X)-m_X,\,0.02\}}{0.02}
 \right),
 \qquad c_{ij}^{(X)}\in\{0,4\},
 \label{eq:q23-integer-cost}
\end{equation}
where $\operatorname{nint}$ is NumPy's \texttt{rint}, which uses ties-to-even.
The value 0.02 is a cap, not a measure of score resolution.  For the
nonnegative gap $d_{ij}=\widehat V_{23}-m_X$, the rounded cost is zero when
$d_{ij}<0.01$.  At the exact half-step, $d_{ij}=0.01$, NumPy's ties-to-even
rule also returns zero.  A unique zero-cost legal state therefore requires
the second-smallest gap to exceed 0.01 strictly.  We retain each raw score
and candidate identifier for reporting outside the integer objective.

With $\lambda=2$, the implemented energy is
\begin{equation}
 E_X(\boldsymbol a,\boldsymbol b)
 = \sum_{i=1}^{52}\sum_{j=1}^{52} c_{ij}^{(X)}a_i b_j
 +2\left(\sum_{i=1}^{52}a_i-1\right)^2
 +2\left(\sum_{j=1}^{52}b_j-1\right)^2.
 \label{eq:q23-accepted-qubo}
\end{equation}
This is a QUBO because $a_i^2=a_i$ and $b_j^2=b_j$
\cite{Glover2019QUBOTutorial,Lucas2014IsingNP}.  In the stored upper-triangular
matrix convention it expands to
\begin{align}
 E_X={}&4-2\sum_i a_i-2\sum_j b_j
 +4\sum_{i<k}a_i a_k+4\sum_{j<\ell}b_j b_\ell
 +\sum_{i,j}c_{ij}^{(X)}a_i b_j .
 \label{eq:q23-expanded-qubo}
\end{align}
The constant offset is therefore 4.  Matrix diagonal entries are $-2$,
within-block couplings are 4, and cross-block entries are the quantized
costs 0 or 4.  For a legal state selecting pair $(i,j)$, the offset and two
linear terms cancel, leaving
\begin{equation}
 E_X(\boldsymbol e_i,\boldsymbol e_j)=c_{ij}^{(X)}.
 \label{eq:q23-legal-energy}
\end{equation}
The offset does not change the optimizer, but it makes the stored legal-state
energy equal to the integer cost and removes ambiguity when matrix and
formula energies are compared.

The two-level construction is designed to preserve the already
enumerated group argmin; it does not preserve all differences between raw
scores.  Y--Ti--$\mathrm{NO}_2$ has a proxy score of $1.491928$, compared
with $1.805747$ for the runner-up Y--Zr--$\mathrm{NO}_2$, giving a
numerical proxy margin of $0.313819$.  This margin is a difference between
predictions, not a confidence interval; the Q2 pool has no physical labels
with which to estimate the error of $\widehat V_{23}$.  The resulting
integer matrix has a single zero-energy legal state because the group
minimum was used to construct the two levels.  Since all cross-block costs
are nonnegative, any violation of either one-hot constraint adds at least
$\lambda=2$ to the energy.  The lowest illegal state is thus two integer
units above the legal minimum.  This separation verifies the encoded
constraint structure, not the physical superiority of the selected
composition.

\subsection{Matrix audit and legal-state enumeration}

The archived Q2 matrix is $104\times104$, has offset 4, and contains
coefficients in $\{-2,0,4\}$.  We enumerated all 2,704 legal one-hot
assignments and recomputed their energies from both
Equation~\eqref{eq:q23-accepted-qubo} and the stored matrix.  The maximum
signed discrepancy was zero.  The unique legal minimum decoded to
candidate~15567, Y--Ti--NO2, with the retained raw score
\begin{equation}
  \widehat V_{23}(\mathrm{Y},\mathrm{Ti},\mathrm{NO}_2)
  =1.4919281724855757.
  \label{eq:q2-best-score}
\end{equation}
The enumerated matrix optimum and the minimum raw score identify the same
candidate.  This check establishes equivalence for the frozen surrogate
objective; it does not establish the unknown physical properties of that
composition.

The corresponding Ising coefficients produced from the accepted QUBO
remain integral after the recorded eight-bit preprocessing, with
$\max|h|=102$ and $\max|J|=1$.  The SDK preprocessing record reports no
loss of the two-unit legal--illegal separation.  These are coefficient and
energy audits of the stored model, not a claim about hardware precision.

\subsection{Kaiwu SDK experiment and Q2 result}

The recorded experiment used Kaiwu SDK~1.4.3 on a local CPU.  It ran
30 fixed random seeds, requesting up to ten returned samples per seed
(300 samples in total).  Simulated annealing used 300 iterations per
temperature, a cooling factor of $\alpha=0.99$, an initial temperature
twice the maximum absolute adapted-matrix entry, and a cutoff temperature
equal to $10^{-3}$ times that entry.  Each seed ran in a single process.
This was a local SDK computation, not an SPQC submission.

The recorded solver time was 20.813~s.  Of the 300 returned samples, 23
($7.67\%$) satisfied both one-hot constraints.  All 23 decoded to
candidate~15567 and came from distinct seeds, so the enumerated optimum
was recovered in 23 of the 30 seed-level attempts with zero regret on the
original proxy score.  The remaining samples failed the one-hot checks.
This legal-sample fraction characterizes this solver configuration; it is
not a probability that the material is physically feasible.

\subsection{Interpretation and limitations}

The fixed-$X$ formulation maps each ordered composition to a
104-variable state and uses integer coefficients compatible with the
recorded eight-bit interface.  Y--Ti--$\mathrm{NO}_2$ is the numerical
argmin of the frozen $\widehat V_{23}$ table and is therefore a priority
candidate for follow-up.  The QUBO experiment confirms the oracle-assisted
encoding of that known proxy minimum; it is neither an independent
discovery nor a measurement of the material's true violation score.

File~11 was also modeled as a separate blind-pool problem over its 529
sparse legal A--B edges, rather than by filtering the full-pool result.
Let $\mathcal L$ denote these edges and define the within-pool normalized
score
\[
 r_{ij}=\frac{\widehat V_{23}(A_i,B_j,\mathrm{NO}_2)-v_{\min}^{\mathcal L}}
 {v_{\max}^{\mathcal L}-v_{\min}^{\mathcal L}},
 \qquad (i,j)\in\mathcal L,
\]
where $v_{\min}^{\mathcal L}$ and $v_{\max}^{\mathcal L}$ are calculated
from the blind-pool predictions alone.  The floating-point model uses
$r_{ij}$ directly; the eight-bit model uses
$c_{ij}^{\mathrm{int}}=4\,\operatorname{round}(4r_{ij})$ for this Q2
pool.  Neither construction uses the enumerated winner identity, and both
matrices were frozen before sampling.  The cross-block coefficient is
set to the scaled cost on $(i,j)\in\mathcal L$ and to $2\lambda=4$
otherwise:
\begin{equation}
 \widetilde c_{ij}=\begin{cases}c_{ij},&(i,j)\in\mathcal L,\\
 2\lambda,&(i,j)\notin\mathcal L.\end{cases}
 \label{eq:blind-sparse-mask}
\end{equation}
The one-hot terms remain, and the decoder additionally requires the
selected edge to belong to $\mathcal L$.  Both local simulated-annealing
runs returned 300/300 legal samples and recovered Y--Ti--NO2 with
$\widehat V_{23}=1.491928$.  The floating-point encoding had one minimum
state and hit it in 25/30 batches; the eight-bit encoding had 11 tied
minimum states and hit the blind-pool minimum in 16/30 batches.  Each
batch is one seed with ten returned reads.  These are oracle-free
blind-pool proxy-search results from local CPU simulated annealing, not
Kaiwu SDK or SPQC-550 results and not physical-label validation.

The final 104-variable Q2 matrix was also submitted to SPQC-550 under the
task ID listed in the appendix.  One of ten returned samples was
legal and decoded to Y--Ti--NO2; the creation-to-finish interval was 22~s
and the platform reported 15.0116~ms evolution time.  This single, small return confirms that the archived matrix produced
a platform response, but it provides no estimate of hardware stability.
The formal solver result remains the Kaiwu SDK~1.4.3 CPU experiment
reported above.
